
\documentclass[12pt]{article}
\usepackage{amsmath}
\usepackage{amsfonts}
\usepackage{amssymb}
\usepackage{geometry}
\geometry{a4paper, margin=0.8in}
\pagenumbering{gobble}

\title{Homework assignment for 02YELMA \\ Chapter: Electric Fields in Matter}
\begin{document}
\date{}
\maketitle
\vspace{-0.8in}

\begin{enumerate}
    
    \item An infinitely long dielectric cylinder of radius $R$ is polarized as $\mathbf{P}(s) = \frac{a}{s}\,\hat{s}$ for $0 < r \le R$, where $a$ is a constant and $s$ is the distance from the axis of the cylinder.

    \begin{enumerate}
        \item Find the bound volume charge density $\rho_b$.
        \item Find the bound surface charge density $\sigma_b$ at $s = R$.
        \item Compute the total bound charge per unit length.
        \item Calculate the total charge and explain why the result is surprising (or not). How can there be a net surface charge if the divergence is zero everywhere?
    \end{enumerate}

    \item A dielectric slab occupies the region $- L \le x \le L$ and is infinite in the $y$ and $z$ directions. The polarization is given by $\mathbf{P}(x) = k x^2 \,\hat{x}$ where $k$ is a constant.

    \begin{enumerate}
        \item Find the bound volume charge density $\rho_b(x)$.
        \item Find the bound surface charge densities at $x = +L$ and $x = -L$.
        \item Compute the total bound charge per unit area.
        \item Determine whether the system is overall neutral. How would you explain this result?
        \item Where is charge concentrated more---near the center or near the boundaries? Explain why.
    \end{enumerate}

    \item A solid sphere of radius $R$ has polarization $\mathbf{P}(r) = k r \,\hat{r}$. 
    \begin{enumerate}
        \item Find $\rho_b$ and $\sigma_b$.
        \item Calculate the electric field inside and outside the sphere. 
        \item Write the expression for the potential $V(0)$ at the center of the sphere.
        \item Compute the volume contribution and the surface contribution. 
        \item Add them and simplify.
        \item Explain why the two contributions do or do not cancel.
    \end{enumerate}

\end{enumerate}

\end{document}
